\xmpheader 6/{A ruled table}
\bigskip
{\offinterlineskip % So the vertical rules are connected.
\def\tablerule{\noalign{\hrule}}% \tablerule constructs a thin rule across the table.
\def\tableskip{\omit&height 9pt&&&\omit\cr}% \tableskip create 9pt of space between entries.
% & separates templates for each column. TeX substitutes the text of the entries for #. We must have
% a strut present in every row of the table; otherwise, the boxes won't butt together properly, 
% and the rules won't join. 
\halign{\tabskip = .7em plus 1 em % glue between columns. Use \vtop for short multiline entries.
\vtop{\hsize=6pc \pretolerance=10000 \hbadness=10000 \normalbaselines \rightskip=3pt plus2em
\noindent \it#\strut}&\vrule #&#\hfil &\vrule #% the rules and midle column.
&\vtop{\hsize=11pc \parindent=0pt \normalbaselineskip=12pt \normalbaselines 
\rightskip=3pt plus2em #}\cr
% The table rows begin here.
\noalign{\hrule height2pt depth2pt \vskip3pt}
\multispan5\bf Some Choice Edible Mushrooms\hfil\strut\cr % The header row spans all the column.
\noalign{\vskip3pt} \tablerule \omit&height 3pt&\omit&&\omit\cr
\bf Botanical&&\bf Common&&\omit\bf Identifying \hfil\cr
\noalign{\vskip -2pt}% close up lines of heading
\bf Name&&\bf Name &&\omit \bf Characteristics \hfil\cr 
\tableskip Pleurotus ostreatus&&Oyster mushroom&&Grow in shelf\kern 1pt like 
% whithout the kern, the `f' and `l' would be too close
clusters on stumps or logs, pink-gray oyster-shaped caps, stem short or absent.\cr 
\tableskip Lactarius hygrophoroides&&Milky hygroph&&Butterscotch-brown cap and stem, 
copious white latex, often on ground in woods near streams.\cr 
\tableskip Morchella esculenta&&White morel&&Conical cap with black pits and white ridges; no gills.  
Often found near old apple trees and dying elms in the spring.\cr
\tableskip Boletus edulus&&King bolete&&Reddish-brown to tan cap with yellow pores 
(white when young), bulbous stem, often near conifers, birch, or~aspen.\cr
\tableskip \tablerule \noalign{\vskip 2pt} \tablerule}
}
